\addcontentsline{toc}{subsubsection}{Solution 67. The Forbidden Coefficient}

\begin{dlsummary}[title={\bfseries Solution 67. The Forbidden Coefficient},breakable]

\begin{proof}
The only solutions are
\[
\boxed{f(x)=x^3\quad\text{and}\quad f(x)=x^3+x.}
\]

\paragraph{1. Remove the cubic contribution.}

Set
\[
g(x)=f(x)-x^3.
\]
The identity
\[
(x+y)^3-x^3-y^3=3xy(x+y)
\]
turns (1) into
\[
0\le g(x+y)-g(x)-g(y)\le(x-y)^2. \tag{4}
\]

Taking \(y=x\) gives
\[
g(2x)=2g(x).
\]
Now apply (4) at \(x/2^n,y/2^n\) and multiply by \(2^n\):
\[
0\le g(x+y)-g(x)-g(y)
\le\frac{(x-y)^2}{2^n}.
\]
Letting \(n\to\infty\), we obtain
\[
\boxed{g(x+y)=g(x)+g(y).} \tag{5}
\]

\paragraph{2. Use positivity to eliminate irregular additive functions.}

For \(0<t<1\), positivity of \(f\) gives
\[
g(t)=f(t)-t^3>-1.
\]
By additivity,
\[
g(t)=g(1)-g(1-t)<g(1)+1.
\]
Thus \(g\) is bounded in absolute value on \((0,1)\); write \(|g(t)|\le M\) there.

For any positive integer \(N\), if \(0<h<1/N\), then
\[
N|g(h)|=|g(Nh)|\le M.
\]
Consequently, \(g(h)\to0\) as \(h\to0^+\). Additivity then gives continuity at every positive input.

Writing \(c=g(1)\), additivity first gives \(g(q)=cq\) for positive rationals, and continuity extends this to every positive real number. Hence
\[
\boxed{f(x)=x^3+cx.} \tag{6}
\]
Positivity for arbitrarily small \(x>0\) forces \(c\ge0\).

\paragraph{3. Discover the forbidden coefficient interval.}

Substitution into (2) gives
\[
f(x)f(y)-f(xy)
=cxy\bigl(c+x^2+y^2-1\bigr)\ge0.
\]
This factorisation was also checked symbolically.[1]

If \(c>0\), division by \(cxy\) yields
\[
c+x^2+y^2\ge1.
\]
Taking arbitrarily small positive \(x,y\), we conclude \(c\ge1\). Therefore
\[
\boxed{c=0\quad\text{or}\quad c\ge1.} \tag{7}
\]

\paragraph{4. Apply the nested-value lock.}

From (6),
\[
f(1)=1+c,
\]
so (3) becomes
\[
(1+c)^3+c(1+c)\le10.
\]
Equivalently,
\[
(c-1)(c^2+5c+9)\le0.
\]
Since \(c\ge0\), the second factor is positive, and thus \(c\le1\).

Together with (7), this leaves exactly
\[
c=0\quad\text{or}\quad c=1.
\]

\paragraph{5. Verify.}

For either candidate, the middle expression in (1) is zero.

For \(f(x)=x^3\), condition (2) is equality and \(f(f(1))=1\).

For \(f(x)=x^3+x\),
\[
f(x)f(y)-f(xy)=xy(x^2+y^2)>0,
\]
and
\[
f(f(1))=f(2)=10.
\]
Both functions satisfy every condition.
\end{proof}

\paragraph{Takeaways.}
\begin{itemize}
\item **Subtract the polynomial suggested by the defect:** This exposes a nearly additive remainder.
\item **Exploit competing scaling rates:** A quadratic error bound cannot survive exact linear scaling.
\item **Use positivity locally:** It supplies the boundedness needed to rule out discontinuous additive behaviour.
\item **Look for coefficient gaps:** Submultiplicativity excludes every \(c\in(0,1)\), rather than merely bounding \(c\).
\end{itemize}

\end{dlsummary}
